\documentclass{article}
\pdfinfoomitdate=1
\pdftrailerid{}

\usepackage{tabularx}
\usepackage{booktabs}

\title{Reflection and Traceability Report on \progname}

\author{\authname}

\date{}

%% Comments

\usepackage{color}

% \newif\ifcomments\commentstrue %displays comments
\newif\ifcomments\commentsfalse %so that comments do not display

\ifcomments
  \newcommand{\authornote}[3]{\textcolor{#1}{[#3 ---#2]}}
  \newcommand{\todo}[1]{\textcolor{red}{[TODO: #1]}}
\else
  \newcommand{\authornote}[3]{}
  \newcommand{\todo}[1]{}
\fi

\newcommand{\wss}[1]{\authornote{magenta}{SS}{#1}}
\newcommand{\plt}[1]{\authornote{cyan}{TPLT}{#1}} %For explanation of the template
\newcommand{\an}[1]{\authornote{cyan}{Author}{#1}}

%% Common Parts

\newcommand{\progname}{RoCam} % PUT YOUR PROGRAM NAME HERE
\newcommand{\authname}{Team \#3, SpaceY
    \\ Zifan Si
    \\ Jianqing Liu
    \\ Mike Chen
    \\ Xiaotian Lou} % AUTHOR NAMES                  

\usepackage{hyperref}
\hypersetup{colorlinks=true, linkcolor=blue, citecolor=blue, filecolor=blue,
    urlcolor=blue, unicode=false}
\urlstyle{same}

\usepackage{indentfirst}
\usepackage{graphicx}

\usepackage{titling}

\pretitle{\begin{center}\includegraphics[width=0.5\textwidth]{../../assets/logo/black.png}\\[0.75em]\LARGE}
        \posttitle{\par\end{center}}

\usepackage[letterpaper, portrait, margin=1in]{geometry}

\usepackage{placeins}

\begin{document}

\maketitle

\plt{Reflection is an important component of getting the full benefits from a
  learning experience.  Besides the intrinsic benefits of reflection, this
  document will be used to help the TAs grade how well your team responded to
  feedback.  Therefore, traceability between Revision 0 and Revision 1 is and
  important part of the reflection exercise.  In addition, several CEAB (Canadian
  Engineering Accreditation Board) Learning Outcomes (LOs) will be assessed based
  on your reflections.}

\section{Changes in Response to Feedback}

\plt{Summarize the changes made over the course of the project in response to
  feedback from TAs, the instructor, teammates, other teams, the project
  supervisor (if present), and from user testers.}

\plt{For those teams with an external supervisor, please highlight how the feedback
  from the supervisor shaped your project.  In particular, you should highlight the
  supervisor's response to your Rev 0 demonstration to them.}

\plt{Version control can make the summary relatively easy, if you used issues
  and meaningful commits.  If you feedback is in an issue, and you responded in
  the issue tracker, you can point to the issue as part of explaining your
  changes.  If addressing the issue required changes to code or documentation, you
  can point to the specific commit that made the changes.  Although the links are
  helpful for the details, you should include a label for each item of feedback so
  that the reader has an idea of what each item is about without the need to click
  on everything to find out.}

\plt{If you were not organized with your commits, traceability between feedback
  and commits will not be feasible to capture after the fact.  You will instead
  need to spend time writing down a summary of the changes made in response to
  each item of feedback.}

\plt{You should address EVERY item of feedback.  A table or itemized list is
  recommended.  You should record every item of feedback, along with the source of
  that feedback and the change you made in response to that feedback.  The
  response can be a change to your documentation, code, or development process.
  The response can also be the reason why no changes were made in response to the
  feedback.  To make this information manageable, you will record the feedback and
  response separately for each deliverable in the sections that follow.}

\plt{If the feedback is general or incomplete, the TA (or instructor) will not
  be able to grade your response to feedback.  In that case your grade on this
  document, and likely the Revision 1 versions of the other documents will be
  low.}

\subsection{SRS and Hazard Analysis}

\subsection{Design and Design Documentation}

\subsection{VnV Plan and Report}

\section{Extras}

\plt{Summarize the extras that were tackled by this project.  Extras
  can include usability testing, code walkthroughs, user documentation, formal
  proof, GenderMag personas, Design Thinking, etc.  Extras should have already
  been approved by the course instructor as included in your problem statement.}

\section{Design Iteration (LO11 (PrototypeIterate))}

\plt{Explain how you arrived at your final design and implementation.  How did
  the design evolve from the first version to the final version?}

\plt{Don't just say what you changed, say why you changed it.  The needs of the
  client should be part of the explanation.  For example, if you made changes in
  response to usability testing, explain what the testing found and what changes
  it led to.}

\section{Design Decisions (LO12)}

\plt{Reflect and justify your design decisions.  How did limitations,
  assumptions, and constraints influence your decisions?  Discuss each of these
  separately.}

\section{Economic Considerations (LO23)}

\plt{Is there a market for your product? What would be involved in marketing your
  product? What is your estimate of the cost to produce a version that you could
  sell?  What would you charge for your product?  How many units would you have to
  sell to make money? If your product isn't something that would be sold, like an
  open source project, how would you go about attracting users?  How many potential
  users currently exist?}

\section{Reflection on Project Management (LO24)}

\plt{This question focuses on processes and tools used for project management.}

\subsection{How Does Your Project Management Compare to Your Development Plan}

\plt{Did you follow your Development plan, with respect to the team meeting plan,
  team communication plan, team member roles and workflow plan.  Did you use the
  technology you planned on using?}

\subsection{What Went Well?}

\plt{What went well for your project management in terms of processes and
  technology?}

\subsection{What Went Wrong?}

\plt{What went wrong in terms of processes and technology?}

\subsection{What Would you Do Differently Next Time?}

\plt{What will you do differently for your next project?}

\section{Ethics and Equity (LO18)}

\plt{Describe how your team applied the principle of equity and universal design
  to ensure equitable treatment of all stakeholders.}

\section{Reflection on Capstone}

\plt{This question focuses on what you learned during the course of the capstone project.}

\subsection{Which Courses Were Relevant}

\plt{Which of the courses you have taken were relevant for the capstone project?}

\subsection{Knowledge/Skills Outside of Courses}

\plt{What skills/knowledge did you need to acquire for your capstone project
  that was outside of the courses you took?}

\end{document}